\documentclass{article}

\usepackage[a4paper, left=1.5cm, right=1.5cm, top=2cm, bottom=2cm]{geometry}

\usepackage{../../../../components/mathtex}

\usepackage{hyperref}
\usepackage{fancyhdr}
\usepackage{ccicons}
\usepackage{caption}
\usepackage{subcaption}
\usepackage{algorithm}
\usepackage{algorithmic}

% Configuration des en-têtes et pieds de page
\pagestyle{fancy}
\fancyhf{} % reset tout

\fancyhead[L]{UFR MATHS INFO}
\fancyhead[C]{Assemblée générale}
\fancyhead[R]{2026-2027}

\fancyfoot[L]{CAI Laurent\\}
\fancyfoot[R]{\thepage}

\begin{document}
\docTitle{Réunion de préparation pour les représentants de maths 2}

\section{Maths en général}
Problème : Pas d'exam, donc dispense, donc problème.\\
Sucer Merle et casser les couilles aux deux autres car c'est le seul qui essaye des choses.\\
On a besoin de note, donc on a BESOIN une solution.
\begin{itemize}
    \item Arrêter le programme à maintenant 
    \item On revoit ce qu'on a fait actuellement
    \item VRAIMENT anonymisé les copies d'examens
    \item Cas où on a un exam ;
    \item Cas où on n'a pas d'exam : 
        \begin{itemize}
            \item Devoir Maison;
            \item QCM de cours;
            \item Exposé;
            \item Kholle;
            \item ...
        \end{itemize}
        Pour les convaincre : 
        \begin{itemize}
            \item 
        \end{itemize}
    \item Cas où on fait la même chose qu'en info : Ne pas le faire le même chose
\end{itemize}
Si jamais les autres profs ne font rien, on menace de bloquer les visios\\

\section{Groupe et actions}
\begin{itemize}
    \item Arrêt : Groupe quotient
\end{itemize}

\section{Intégration et probabilité}


\section{Calcul différentielle}
\begin{itemize}
    \item Arrêt : S'arrêter à la différentiabilité;
\end{itemize}

\section{Anglais}
Avis : Dispense
Cas pour le groupe des MI de jeudi : 
\end{document}